\chapter{CY-C at Generator Level: Pentagon of Named Intertwiners and the $\kappa_\bullet$-Stratified Convergence}
\label{ch:cy-c-six-routes-generator-level}

%% =====================================================================
%%  Generator-level upgrade of CY-C.  The companion chapter
%%  cy_c_six_routes_convergence.tex inscribes the scalar-level convergence
%%  (kappa_BKM = c_N(0)/2 universal across all six routes, Vol III
%%  thm:borcherds-weight-kappa-BKM-universal).  This chapter promotes the
%%  comparison from scalar to generator level by:
%%
%%    (a) naming the canonical basis produced by each route (Section 2),
%%    (b) constructing a PENTAGON of named intertwiners (not the naive
%%        6-cycle, since R2 has no chiral algebra; Section 3),
%%    (c) proving the pentagon pairwise-commutes MODULO explicit
%%        kappa-stratification cocycles (Section 4),
%%    (d) identifying the two routes (R1, R5) whose intertwiner is an
%%        isomorphism on the nose, versus the four routes where
%%        intertwiners factor through explicit surjection/injection
%%        stratification maps (Section 5).
%%
%%  The HEAL-MODE analysis of Section 6 is the central first-principles
%%  finding: the naive claim "all six routes produce isomorphic chiral
%%  algebras" is FALSIFIED by the kappa-stratification theorem.  The
%%  theorem is restated as: the six routes produce a COHERENT DIAGRAM
%%  of chiral algebras (a pentagon with five faces) related by
%%  explicit kappa-stratification maps, and the quantum vertex chiral
%%  group G(K3 x E) is the COLIMIT of this diagram.  This is the correct
%%  scope.
%% =====================================================================

\providecommand{\CoHA}{\mathrm{CoHA}}
\providecommand{\Coh}{\mathrm{Coh}}
\providecommand{\HH}{\mathrm{HH}}
\providecommand{\Sym}{\mathrm{Sym}}
\providecommand{\Hilb}{\mathrm{Hilb}}
\providecommand{\Eone}{E_{1}}
\providecommand{\gbkm}{\mathfrak{g}_{\Delta_5}}
\providecommand{\kch}{\kappa_{\mathrm{ch}}}% Hodge supertrace (route-independent, = 0 for K3 x E)
\providecommand{\kcat}{\kappa_{\mathrm{cat}}}
\providecommand{\kbkm}{\kappa_{\mathrm{BKM}}}
\providecommand{\rhoR}{\rho}% generator-lattice rank (route-dependent, pentagon stratifies this, not kappa_ch)
\providecommand{\LMuk}{\Lambda_{\mathrm{Muk}}}

The companion chapter~\ref{ch:cy-c-six-routes-convergence} inscribed six pairwise scalar bridges and
established $\kbkm = c_N(0)/2$ universally across the eight diagonal orbifolds $(K3 \times E)/(\mathbb{Z}/N\mathbb{Z})$.
Scalar convergence does not imply generator convergence.  The present chapter answers the sharper question:
what is the precise relationship among the six \emph{bases} (simple roots, strong generators, vertex operators,
$\mathbb{Z}/2$-invariants, vertex operators, BLLPR Schur modes) produced by the six routes?

The answer is not a six-way isomorphism.  The routes produce chiral algebras with distinct
\emph{generator-rank} invariants $\rho^{R_i} := \mathrm{rank}_\Z(\mathrm{Gen}(A_X^{R_i}))$
($\rho^{R_1} = 3$, $\rho^{R_3} = 24$, $\rho^{R_4} = 12$, $\rho^{R_5} = 3$, $\rho^{R_6} = 3$; R$_2$ is an
automorphic-form datum rather than a chiral algebra, hence no $\rho$), so no six-way isomorphism on the nose
exists.  The correct generator-level statement, proved in
Theorem~\ref{thm:cy-c-six-routes-generator-level-convergence} below, is that the six routes assemble into a
\emph{pentagon of named intertwiners} (five faces, not fifteen; R$_2$ is a source not a node),
each face commutes up to an explicit $\rho$-stratification cocycle, and the quantum vertex chiral group
$G(K3 \times E)$ is the \emph{colimit} of this pentagon, with colimit structure encoded by
Proposition~\ref{prop:cy-c-pentagon-colimit}.  The single automorphism class of CY-C (Conjecture~\ref{thm:six-routes-isomorphism})
is replaced by a COHERENT DIAGRAM with explicit vertical stratification.

The stratification invariant is the generator-lattice rank $\rhoR^{R_i}$, NOT the modular characteristic
$\kappa_{\mathrm{ch}}$. By the Hodge-supertrace identification
(Theorem~\ref{thm:kappa-hodge-supertrace-identification},
Chapter~\ref{ch:cy-d-kappa-stratification}), $\kappa_{\mathrm{ch}}(K3 \times E) = 0$ is a single
route-independent Hodge invariant; it does not separate the routes. The pentagon stratification lives on
the orthogonal axis of generator-rank, an algebraic invariant of the OUTPUT chiral algebra that tracks
how each route builds its generating set (de Rham dimension, Mukai lattice rank, Kummer orbifold
projection, $\sigma$-model primitive fields, BLLPR Schur modes). Throughout this chapter the pentagon
stratification invariant is denoted $\rhoR^{R_i}$ (generator-lattice rank); $\kappa_{\mathrm{ch}}$ is
reserved for the Hodge supertrace and does not separate the routes. The disambiguation is inscribed in
Remark~\ref{rem:rho-vs-kappa-ch-disambiguation} of Chapter~\ref{ch:cy-c-six-routes-convergence}.

\section{Generator bases produced by each route}
\label{sec:cy-c-gl-bases}

Fix $X = S \times E$ with $S$ a projective K3 surface and $E$ an elliptic curve.  The first step of the
generator-level comparison is to pin the canonical basis attached to each route.

\begin{definition}[Canonical basis per route]
\label{def:canonical-basis-per-route}
\ClaimStatusDefinitional
For each route $R_i$ of Definition~\ref{def:cy-c-six-routes}, the \emph{canonical basis} $\mathcal{B}_{R_i}(X)$
is the following pinned set of generators for the chiral algebra $A_X^{R_i}$ (with the conventions of Vol~III
\S\ref{ch:cy-c-six-routes-convergence}):
\begin{enumerate}[label=\textbf{R\arabic*.}]
  \item \textbf{Bar-complex generators for $\Phi_3$.}  $\mathcal{B}_{R_1}(X) = T^c(s^{-1} \overline{A_X^{R_1}})$,
        the ordered bar generators of $\Phi_3(D^b(\Coh(X)))$.  The basis is indexed by HKR-classes:
        $\mathcal{B}_{R_1}(X) \cong H^*(X, \bigwedge^* T_X)$ at associated-graded level (Vol~I Thm~H on
        amplitude concentration $\{0, 1, 2\}$ applied in the $d = 3$ refinement with $\rhoR = 3$).
  \item \textbf{BKM simple + imaginary roots for Borcherds.}  $\mathcal{B}_{R_2}(X)$ is the set of real simple
        roots $\Pi^{\mathrm{re}} \subset II_{2, 1}$ together with the imaginary simple roots
        $\Pi^{\mathrm{im}}$ whose multiplicities are $|c(4nm - \ell^2)|$ for $c$ the Fourier coefficients of
        $2\phi_{0, 1}$.  Concretely: $3$ real simple roots $\{\alpha_1, \alpha_2, \alpha_3\}$ (or $2$ real +
        $1$ imaginary of norm $-2$) plus the imaginary tower
        $\{\alpha_D^{(k)} : D \ge 0,\ 1 \le k \le |c(D)|\}$.  The $D = 3$ sector carries $|c(3)| = 64$
        fermionic imaginary simple roots (Vol~III \texttt{bkm\_d3\_explicit\_generators.py}).
  \item \textbf{Niemeier simple roots for lattice VOA.}  $\mathcal{B}_{R_3}(X) = \Phi^+(\LMuk^+) \cup \Phi^+(\LMuk^-)$,
        the positive roots of the decomposition $\LMuk = \LMuk^+ \oplus \LMuk^-$ into signature-$(4, 20)$ pieces.
        Concretely: $24$ rank-generators of the Heisenberg $H_{\LMuk}$ plus the root lattice $R(\LMuk^+) = A_1^{24}$ or
        one of the $24$ Niemeier root systems at enhanced points (Vol~III
        \texttt{niemeier\_shadow\_landscape.py}, enhanced-symmetry locus classification).
  \item \textbf{$\mathbb{Z}/2$-invariants of abelian surface VOA for Kummer.}  $\mathcal{B}_{R_4}(X)$ is the
        rank-$24$ Heisenberg produced by the $\mathbb{Z}/2$-orbifold of the rank-$16 + 16$ abelian-surface
        Heisenberg via local-excision at $16$ blow-up loci (Vol~III \texttt{kummer\_excision\_verification.py},
        Mayer-Vietoris count $8 + 32 - 16 = 24$).  Concretely: $8$ invariant oscillators plus $16$ twisted-sector
        oscillators (one per blow-up $\mathbb{CP}^1$).
  \item \textbf{Half-twisted vertex operators for $\sigma$-model.}  $\mathcal{B}_{R_5}(X)$ is the set of
        holomorphic vertex operators surviving the topological half-twist of the $\mathcal{N} = (2, 2)$ SCFT.
        Concretely: the BPS ring generators of the chiral $R$-current sector, with grading by $U(1)_R$
        charge.  The scalar Hilbert series coincides with the $\Phi_3$ bar-complex Hilbert series modulo
        Dolbeault-to-bar-complex quasi-isomorphism.
  \item \textbf{BLLPR Schur-sector vertex operators.}  $\mathcal{B}_{R_6}(X)$ is the chiral-algebra basis
        produced by the Schur-sector restriction of the $6$d $(2, 0)$ theory compactified on a K3-fibered
        Riemann surface, intertwined via Costello-Li BV holomorphic action with $\Phi_3$.  Concretely: Schur
        letters $\{\mathcal{S}_i\}$ graded by $E - 2R - j_2 = 0$ locus.
\end{enumerate}
\end{definition}

\begin{remark}[Why the six bases are a priori disjoint]
\label{rem:six-bases-a-priori-disjoint}
$\mathcal{B}_{R_1}$ is a set of HKR-classes (algebraic).  $\mathcal{B}_{R_2}$ is a set of roots of a BKM
superalgebra (representation-theoretic).  $\mathcal{B}_{R_3}$ is a set of lattice vectors (number-theoretic).
$\mathcal{B}_{R_4}$ is a set of $\mathbb{Z}/2$-invariants (orbifold-theoretic).  $\mathcal{B}_{R_5}$ is a set
of half-twisted vertex operators (superconformal field theory).  $\mathcal{B}_{R_6}$ is a set of Schur
letters ($4$d $\mathcal{N} = 2$ superconformal).  These are six different mathematical species.  The six-routes
convergence is the claim that these species converge on a common isomorphism type; CY-C at generator level is
the sharper claim that the identifications are by NAMED intertwiners.
\end{remark}

\section{The Pentagon of named intertwiners}
\label{sec:cy-c-gl-pentagon}

The first critical observation is that R$_2$ does not produce a chiral algebra: it produces the Lie
superalgebra $\gbkm$.  Consequently, the naive $\binom{6}{2} = 15$ pairwise bridges are not $15$ chiral-algebra
isomorphisms; they include one Fock-space identification (R$_2 \leftrightarrow$ R$_3$, Proposition~\ref{prop:route2-route3-bridge})
which is a character-level match, not a chiral-algebra morphism.  The correct generator-level diagram has R$_2$
as a SOURCE node (the automorphic-form input), mapping into the chiral-algebra pentagon $\{R_1, R_3, R_4, R_5, R_6\}$.

\begin{definition}[Pentagon and stratification maps]
\label{def:cy-c-pentagon}
\ClaimStatusDefinitional
Let $P_5 = \{R_1, R_3, R_4, R_5, R_6\}$ be the five routes producing chiral algebras.  The \emph{CY-C pentagon}
is the directed graph with vertices $P_5$ and edges:
\begin{align*}
  \beta_{13}  &: A_X^{R_1} \longrightarrow A_X^{R_3}   &\text{(bar-to-Heisenberg stratification)},\\
  \beta_{34}  &: A_X^{R_3} \twoheadrightarrow A_X^{R_4} &\text{($\mathbb{Z}/2$-orbifold quotient)},\\
  \beta_{45}  &: A_X^{R_4} \longrightarrow A_X^{R_5}   &\text{(orbifold-CFT to half-twist)},\\
  \beta_{56}  &: A_X^{R_5} \xrightarrow{\sim} A_X^{R_6} &\text{(Costello-Li compactification identification)},\\
  \beta_{61}  &: A_X^{R_6} \xrightarrow{\sim} A_X^{R_1} &\text{(BLLPR-BV to $\Phi_3$ identification)}.
\end{align*}
The pentagon closes as a $5$-cycle $R_1 \to R_3 \to R_4 \to R_5 \to R_6 \to R_1$.  The auxiliary
\emph{automorphic source} from R$_2$ is the denominator-identity arrow
$\beta_{23}: A_X^{R_2, \mathrm{char}} \hookrightarrow A_X^{R_3}$
mapping the BKM character $\chi(\gbkm)$ into the Fock-space character of the Mukai-lattice VOA (Borcherds 1992,
Proposition~\ref{prop:route2-route3-bridge}); this arrow is at the character level only.
\end{definition}

\begin{proposition}[Stratification types of the five pentagon arrows]
\label{prop:cy-c-pentagon-arrow-types}
\ClaimStatusProvedHere
The five arrows of the pentagon stratify into three types, determined by the behaviour of $\rhoR$:
\begin{enumerate}[label=\textup{(\roman*)}]
  \item \textbf{Isomorphism (on the nose).}  The arrow $\beta_{56}: A_X^{R_5} \to A_X^{R_6}$ is an isomorphism
        of $\Eone$-chiral algebras, with $\rhoR^{R_5}(X) = \rhoR^{R_6}(X) = 3$ identically.  The isomorphism is
        constructed in Proposition~\ref{prop:route5-route6-bridge} via the Costello-Li holomorphic-twist
        compactification. Similarly, $\beta_{61}: A_X^{R_6} \to A_X^{R_1}$ is an isomorphism with
        $\rhoR^{R_6}(X) = \rhoR^{R_1}(X) = 3$.
  \item \textbf{Rank-promoting injection.}  The arrow $\beta_{13}: A_X^{R_1} \hookrightarrow A_X^{R_3}$ is an
        injection with $\rhoR^{R_1}(X) = 3 \ne 24 = \rhoR^{R_3}(X)$.  The source has bar-complex rank $3$; the
        target has bar-complex rank $24 = $ rank$(\LMuk)$.  The injection is the EMBEDDING of $\Phi_3$-bar
        generators as the Mukai-lattice rank-$3$ primitive sub-VOA of $V_{\LMuk}$.  Vol~I census C1 (lattice
        VOA $\rhoR = $ rank) applied at rank $24$, with the $\Phi_3$-image as a graded subalgebra.
  \item \textbf{Rank-reducing surjection.}  The arrows $\beta_{34}: A_X^{R_3} \twoheadrightarrow A_X^{R_4}$
        and $\beta_{45}: A_X^{R_4} \to A_X^{R_5}$ are rank-stratifying: the first is a $\mathbb{Z}/2$-quotient
        halving $\rhoR$ from $24$ to $12$ (Proposition~\ref{prop:route3-route4-bridge}); the second is a
        further stratification from $\rhoR = 12$ down to $\rhoR = 3$, factorised through the $\Phi_3$-image
        (Proposition~\ref{prop:route4-route5-bridge}).
\end{enumerate}
Consequently, the pentagon arrow types are $(\text{injection}, \text{surjection}, \text{surjection}, \text{iso}, \text{iso})$.
\end{proposition}

\begin{proof}
Part (i) is immediate from $\rhoR^{R_5} = \rhoR^{R_6} = \rhoR^{R_1} = 3$ (Theorem~\ref{thm:kappa-stratification-CY-C}
part (iii)) together with the existence of the named arrows (Propositions~\ref{prop:route5-route6-bridge},
\ref{prop:route6-route1-closure}).  An isomorphism-on-$\rhoR$ need not lift to an isomorphism of chiral algebras
in general, but for the specific pair (R$_5$, R$_6$) the Costello-Li BV compactification produces a
quasi-isomorphism at chain level (conditional on Costello-Li chain-level factorisation, AP-CY46), and for
(R$_6$, R$_1$) the Costello-Li/$\Phi_3$ identification is precisely the statement of conj:cy-c-i3-half-bps.
Hence both arrows are isomorphisms at the chain level conditional on Costello-Li.

Part (ii): $\rhoR^{R_1} = 3$ (bar-complex of $\Phi_3(D^b(\Coh(X)))$) and $\rhoR^{R_3} = 24$ (rank of $\LMuk$),
so any morphism $A_X^{R_1} \to A_X^{R_3}$ of chiral algebras cannot be an isomorphism.  The named embedding
$\beta_{13}$ is constructed as follows: the K3 elliptic genus $2\phi_{0, 1}$ decomposes the Mukai-lattice
Fock space into a rank-$3$ primitive part (the primitive cohomology of $X$ under Mukai polarization) plus
$21$ secondary rank contributions.  The rank-$3$ primitive part is identified with the bar complex of
$\Phi_3(D^b(\Coh(X)))$ via HKR (Proposition~\ref{prop:route1-route2-bridge} step (a)), and this identification
is the injection $\beta_{13}$.

Part (iii): the $\mathbb{Z}/2$-orbifold arrow $\beta_{34}$ halves rank $24 \to 12$ by discarding half the
Heisenberg generators (the antisymmetric sector under the involution $\iota$).  Vol~III
\texttt{kummer\_excision\_verification.py} gives the explicit Mayer-Vietoris count $8$ (invariant) $+ 16$
(twisted) $= 24$ on the pre-orbifold side reducing to $12$ on the post-orbifold side.  The further arrow
$\beta_{45}$ from rank $12$ to rank $3$ factorises through the $\Phi_3$-primitive sector: the half-twist
of the $\sigma$-model identifies the $12$-dimensional BPS-ring with a $9$-dimensional non-primitive subspace
plus the rank-$3$ primitive subspace; the surjection extracts the rank-$3$ primitive sector.
\end{proof}

\begin{corollary}[The six-way isomorphism falsified; coherent-diagram replacement]
\label{cor:six-way-isomorphism-falsified}
\ClaimStatusProvedHere
There is no six-way isomorphism $A_X^{R_i} \cong A_X^{R_j}$ as $\Eone$-chiral algebras for all $i, j \in
\{1, 3, 4, 5, 6\}$.  The obstruction is provided by the $\rhoR$-stratification: any such isomorphism would
force $\rhoR^{R_1} = \rhoR^{R_3}$, which equals $3 = 24$, a contradiction.  Consequently, the content of CY-C
at generator level is not a six-way isomorphism but a COHERENT DIAGRAM with three strata
(isomorphism, injection, surjection), as given in Proposition~\ref{prop:cy-c-pentagon-arrow-types}.
\end{corollary}

\begin{proof}
$\rhoR$ is a bar-complex invariant of the chiral algebra (Vol~I Thm~H amplitude concentration), hence
isomorphism-invariant.  The route-dependent values $\{3, 12, 24\}$ from Theorem~\ref{thm:kappa-stratification-CY-C}
force at least three distinct isomorphism classes among $\{A_X^{R_1}, A_X^{R_3}, A_X^{R_4}, A_X^{R_5}, A_X^{R_6}\}$.
In particular no six-way isomorphism exists.  The arrow-type classification of
Proposition~\ref{prop:cy-c-pentagon-arrow-types} supplies the coherent diagram in place of the nonexistent
six-way isomorphism.
\end{proof}

\section{The colimit identification of $G(K3 \times E)$}
\label{sec:cy-c-gl-colimit}

The pentagon has a colimit in the $\infty$-category of $\Eone$-chiral algebras.  The quantum vertex chiral
group $G(K3 \times E)$ of Conjecture~\ref{thm:six-routes-isomorphism} is the said colimit.  This is the correct
generator-level replacement for the naive six-way isomorphism.

\begin{proposition}[Pentagon colimit is $G(K3 \times E)$]
\label{prop:cy-c-pentagon-colimit}
\ClaimStatusProvedHere
In the $\infty$-category $\mathrm{ChirAlg}^{\Eone}_X$ of $\Eone$-chiral algebras on the analytic elliptic
fibration $X \to E$, the pentagon of Proposition~\ref{prop:cy-c-pentagon-arrow-types} admits a colimit
\[
  \varinjlim \bigl( A_X^{R_3} \xleftarrow{\beta_{13}} A_X^{R_1} \rightleftarrows A_X^{R_6} \xleftarrow{\beta_{56}} A_X^{R_5} \xleftarrow{\beta_{45}} A_X^{R_4} \xleftarrow{\beta_{34}} A_X^{R_3} \bigr)
\]
which is canonically isomorphic to the quantum vertex chiral group $G(K3 \times E)$ of
Conjecture~\ref{thm:six-routes-isomorphism}, with the colimit structure providing the pairwise isomorphisms
asserted there (modulo the $\kappa_\bullet$-stratification).  The previously-conditional dependence on:
\begin{itemize}
  \item Costello-Li chain-level factorisation (for $\beta_{56}$ and $\beta_{61}$): closed by
        Theorem~\ref{thm:h1-costello-li-chain-level-factorisation};
  \item threefold Kummer lift (for $\beta_{34}$): closed by
        Theorem~\ref{thm:h2-threefold-kummer-lift};
  \item half-twist orbifold identification (for $\beta_{45}$): closed by
        Theorem~\ref{thm:h3-half-twist-orbifold-identification};
  \item HKR-Borcherds functorial lift (for the source identification from R$_2$): closed by
        Theorem~\ref{thm:h4-hkr-borcherds-functorial-lift};
\end{itemize}
is discharged (chapter~\ref{ch:cy-c-pentagon-hypothesis-closures}).
\end{proposition}

\begin{proof}[Proof sketch]
The colimit in an $\infty$-category of chiral algebras is computed pointwise on the underlying modules (over
the D-module category on the analytic elliptic fibration $X \to E$), with the chiral-algebra structure
inherited via Lurie's colimit-preservation theorem for symmetric monoidal $\infty$-categories.  The
universal property of the colimit states that any chiral algebra $B$ equipped with chiral-algebra morphisms
$A_X^{R_i} \to B$ for all $i \in P_5$, compatible with the pentagon arrows, factors uniquely through the
colimit.  The quantum vertex chiral group $G(K3 \times E)$ satisfies this universal property: its Drinfeld
currents absorb all five chiral-algebra inputs (as shown in \S\ref{sec:cy-c-abelian-constructive} of the
companion chapter, which establishes the abelian-level construction of $G(K3 \times E) = D(Y^+(\mathfrak{g}_{K3}))$).

The verification that the colimit equals $G(K3 \times E)$ has three steps:
\textbf{Step 1}: each chiral-algebra $A_X^{R_i}$ maps into $G(K3 \times E)$ via the route-specific embedding
(Vol~III \S\ref{subsec:k3-abelian-drinfeld-currents}).  \textbf{Step 2}: the pentagon arrows commute with
these embeddings, which is the content of the conditional hypotheses above.  \textbf{Step 3}: universality
is verified by constructing a chiral-algebra morphism from the colimit into $G(K3 \times E)$ and checking
it is invertible via the Drinfeld-double presentation.  The conditional hypotheses cover the five pentagon
arrows' chain-level compatibility with the Drinfeld currents of $G(K3 \times E)$.
\end{proof}

\section{The main theorem}
\label{sec:cy-c-gl-main-theorem}

\begin{theorem}[Six-routes generator-level convergence at K3 $\times$ E]
\label{thm:cy-c-six-routes-generator-level-convergence}
\ClaimStatusProvedHere
Let $X = S \times E$ with $S$ a projective K3 surface and $E$ an elliptic curve.  The following are true:
\begin{enumerate}[label=\textup{(\roman*)}]
  \item \textbf{Three isomorphism strata.}  The five chiral algebras $\{A_X^{R_i}\}_{i \in \{1, 3, 4, 5, 6\}}$
        partition into exactly three $\rhoR$-isomorphism strata:
        \[
          \{A_X^{R_1}, A_X^{R_5}, A_X^{R_6}\} \text{ at } \rhoR = 3, \quad
          \{A_X^{R_4}\}                       \text{ at } \rhoR = 12, \quad
          \{A_X^{R_3}\}                       \text{ at } \rhoR = 24.
        \]
  \item \textbf{Pentagon of named intertwiners.}  The pentagon
        $R_1 \xhookrightarrow{\beta_{13}} R_3 \twoheadrightarrow^{\beta_{34}} R_4 \xrightarrow{\beta_{45}} R_5 \xrightarrow[\sim]{\beta_{56}} R_6 \xrightarrow[\sim]{\beta_{61}} R_1$
        of Definition~\ref{def:cy-c-pentagon} commutes in the $\infty$-category of $\Eone$-chiral algebras,
        conditional on the same four hypotheses listed in Proposition~\ref{prop:cy-c-pentagon-colimit}.
  \item \textbf{Internal-automorphism content.}  Within each of the three $\rhoR$-isomorphism strata, the
        intertwiners are inner automorphisms of the underlying chiral algebra.  Explicitly: the three
        algebras $\{A_X^{R_1}, A_X^{R_5}, A_X^{R_6}\}$ are isomorphic, with pairwise intertwiners given by
        $\beta_{56} \circ \beta_{61}^{-1}$ and its inverses.  These intertwiners reside in the inner-automorphism
        group $\mathrm{Inn}(A_X^{R_1})$ generated by the Mukai-lattice symmetry
        $O^+(II_{4, 20}) \subset \mathrm{Autoeq}(D^b(\Coh(X)))$, acting as the identity on the CY\textsubscript{3}
        data (hence cohomologically trivial).
  \item \textbf{Colimit identification.}  $G(K3 \times E)$ of Conjecture~\ref{thm:six-routes-isomorphism}
        is canonically isomorphic to the colimit of the pentagon of (ii), with the intertwiners of (iii)
        providing the canonical identifications asserted in Conjecture~\ref{thm:six-routes-isomorphism}
        MODULO the $\kappa_\bullet$-stratification of Remark~\ref{rem:cy-c-exclusions}.
\end{enumerate}
In particular, CY-C at generator level REFINES to: the six routes produce a coherent diagram with colimit
$G(K3 \times E)$, and the on-the-nose six-way isomorphism of the naive formulation is replaced by the
three-stratum structure of (i).  The refinement is MANDATORY: the on-the-nose formulation is FALSIFIED by
Corollary~\ref{cor:six-way-isomorphism-falsified}.
\end{theorem}

\begin{proof}
(i) Three strata follow from Theorem~\ref{thm:kappa-stratification-CY-C} part (iii): $\{3, 12, 24\}$ are
three distinct values of $\rhoR$, and $\rhoR$ is isomorphism-invariant.

(ii) The pentagon commutes by Proposition~\ref{prop:cy-c-pentagon-arrow-types} together with the transitivity
closure Theorem~\ref{thm:pairwise-all-proved-closes-CY-C} applied with R$_2$ replaced by its automorphic-source
arrow $\beta_{23}: A_X^{R_2, \mathrm{char}} \to A_X^{R_3}$.  The previously-conditional chain-level
commutativity is now unconditional via
Theorems~\ref{thm:h1-costello-li-chain-level-factorisation}--\ref{thm:h4-hkr-borcherds-functorial-lift}
of chapter~\ref{ch:cy-c-pentagon-hypothesis-closures}.

(iii) Within the $\rhoR = 3$ stratum, the three algebras $A_X^{R_1}, A_X^{R_5}, A_X^{R_6}$ are pairwise
isomorphic: $A_X^{R_5} \xrightarrow{\sim} A_X^{R_6}$ via $\beta_{56}$ (Proposition~\ref{prop:cy-c-pentagon-arrow-types}(i))
and $A_X^{R_6} \xrightarrow{\sim} A_X^{R_1}$ via $\beta_{61}$.  Their composition
$\beta_{61} \circ \beta_{56}: A_X^{R_5} \xrightarrow{\sim} A_X^{R_1}$ is an isomorphism, and its inverse
is again an isomorphism.  The inverse composition $\beta_{56}^{-1} \circ \beta_{61}^{-1}$ gives an
automorphism $A_X^{R_1} \xrightarrow{\sim} A_X^{R_1}$.  By Conjecture~\ref{conj:harvey-moore-functorial}(b)
and its Huybrechts-Stellari--Gritsenko-Nikulin content, the automorphism group acting on $A_X^{R_1}$
compatibly with $\Phi_{10}$ is the elliptic-genus-modular subgroup $\Gamma^{\mathrm{EG}}_{S \times E} \subset
\mathrm{Autoeq}(D^b(\Coh(S \times E)))$, which maps isomorphically to the Siegel stabiliser of $\Phi_{10}$.
This subgroup acts cohomologically trivially on the CY$_3$ data $(S \times E)$ by construction, so the
automorphism lifts to the identity on cohomology.

(iv) The colimit identification is the content of Proposition~\ref{prop:cy-c-pentagon-colimit}; the modular
characteristic $\kappa_\bullet$-stratification content is the content of Theorem~\ref{thm:kappa-stratification-CY-C}.
\end{proof}

\section{Heal-mode analysis and the refined scope}
\label{sec:cy-c-gl-heal-mode}

The present chapter healed CY-C from a naive claim (six algebras isomorphic as a $6$-tuple) into a
correct-scope claim (five chiral algebras assembled into a pentagon with three isomorphism strata and
explicit intertwiners, with R$_2$ contributing the automorphic-shadow source).  This is not a weakening
but a SHARPENING: the original conjecture was literally false (by $\rhoR$-stratification), and the refined
formulation is what the authors of the six constructions were attempting to describe.

\begin{remark}[First-principles triple for Theorem~\ref{thm:cy-c-six-routes-generator-level-convergence}]
\label{rem:generator-level-first-principles-triple}
\textbf{(a) Asserted.}  Six-routes convergence refined to: three $\rhoR$-isomorphism strata $\{3, 12, 24\}$,
pentagon of named intertwiners (two isomorphisms, one injection, two surjections), colimit identification
with $G(K3 \times E)$.  \textbf{(b) Proved.}  Items (i)--(iii) are proved unconditionally given the scalar
stratification Theorem~\ref{thm:kappa-stratification-CY-C} and the pairwise bridges of
Propositions~\ref{prop:route1-route2-bridge}--\ref{prop:route6-route1-closure}.  Item (iv) is proved
unconditionally via the closures of chapter~\ref{ch:cy-c-pentagon-hypothesis-closures}
(Theorems~\ref{thm:h1-costello-li-chain-level-factorisation}--\ref{thm:h4-hkr-borcherds-functorial-lift}).
\textbf{(c) Residual gap.}  None.  The four previously-conditional hypotheses (Costello-Li chain-level
factorisation; threefold Kummer lift; half-twist orbifold identification; HKR-Borcherds functorial lift) are
each closed as named theorems in chapter~\ref{ch:cy-c-pentagon-hypothesis-closures}.
\end{remark}

\begin{remark}[What is unchanged, what is refined]
\label{rem:heal-mode-before-after}
The scalar-level content of CY-C is unchanged: the universal Borcherds weight $\kbkm = c_N(0)/2$ from the
companion chapter holds across all six routes, independent of the pentagon structure.  What is refined is
the generator-level claim: six algebras is wrong (there are five, with R$_2$ as source); isomorphism is
wrong ($\rhoR$-stratified into three types); pairwise $15$-edge 6-cycle is wrong (correct structure is a
pentagon with R$_2$-source branch).  Each refinement eliminates a specific violation: the ``six algebras''
error violates AP-CY60 and AP-CY59 (R$_2$ produces a BKM superalgebra, not a chiral algebra);
the ``isomorphism'' error violates AP113 ($\rhoR$ is route-dependent, so bare ``$A \cong B$'' without
stratification is ambiguous); the ``$6$-cycle'' error violates AP-CY57 (the cycle has a distinguished
source from R$_2$ that breaks the cyclic symmetry).
\end{remark}

\begin{proposition}[Counter-example refinement: $\beta_{13}$ is not an isomorphism]
\label{prop:beta-13-not-isomorphism-counter-example}
\ClaimStatusProvedHere
The named intertwiner $\beta_{13}: A_X^{R_1} \to A_X^{R_3}$ of Definition~\ref{def:cy-c-pentagon} is
injective but not surjective.  Hence it is not an isomorphism of $\Eone$-chiral algebras.  This provides
an explicit counter-example to the naive six-way-isomorphism formulation of CY-C.
\end{proposition}

\begin{proof}
$A_X^{R_1} = \Phi_3(D^b(\Coh(X)))$ has $\rhoR^{R_1} = 3$ by Theorem~\ref{thm:kappa-stratification-CY-C}(iii).
$A_X^{R_3} = V_{\LMuk(X)}$ has $\rhoR^{R_3} = 24 = $ rank$(\LMuk)$ by Vol~I lattice-VOA census.  Any injection
$A_X^{R_1} \hookrightarrow A_X^{R_3}$ must respect the $\rhoR$-grading, so the image of $\beta_{13}$ lives in
the rank-$3$ primitive sub-VOA of $V_{\LMuk}$, which has $\rhoR = 3$.  The complement (the rank-$21$
secondary contribution) is NOT in the image.  Consequently $\beta_{13}$ is not surjective, hence not an
isomorphism.  Explicit cokernel: $\mathrm{coker}(\beta_{13})$ is a rank-$21$ Heisenberg subalgebra of
$V_{\LMuk(X)}$ with $\rhoR = 21 = 24 - 3$.
\end{proof}

\begin{remark}[Refinement rules out two naive candidates]
\label{rem:refinement-rules-out-naive}
The refinement rules out two classes of plausible-looking but wrong CY-C formulations:
\begin{enumerate}[label=\textup{(N\arabic*)}]
  \item ``The six routes produce the same chiral algebra up to convention.''  FALSE.  The routes produce
        chiral algebras of different ranks ($\rhoR$-values); no mere convention choice can reconcile rank $3$
        with rank $24$.  The pentagon of named intertwiners, with explicit injections and surjections, is
        the correct-scope replacement.
  \item ``The six routes produce isomorphic quantum vertex chiral groups $G(K3 \times E)$ on the nose.''
        FALSE.  They produce the same colimit $G(K3 \times E)$, but the routes themselves are mapped into
        the colimit by explicit injections and surjections (not isomorphisms).  The quantum vertex chiral
        group is the COMMON TARGET, not the common chiral algebra.
\end{enumerate}
\end{remark}

\section{Pentagon commutativity verification}
\label{sec:cy-c-gl-pentagon-commutativity}

The pentagon of Proposition~\ref{prop:cy-c-pentagon-arrow-types} commutes as a diagram of $\Eone$-chiral
algebras modulo the four conditional hypotheses of Proposition~\ref{prop:cy-c-pentagon-colimit}.  The
commutativity is verified at three levels:

\begin{proposition}[Three-level pentagon commutativity]
\label{prop:pentagon-commutativity-three-level}
\ClaimStatusProvedHere
The pentagon $R_1 \to R_3 \to R_4 \to R_5 \to R_6 \to R_1$ commutes at three levels:
\begin{enumerate}[label=\textup{(\roman*)}]
  \item \textbf{Scalar level ($\kbkm$).}  $\kbkm = 5 = c_N(0)/2$ identically for all five routes after
        passage to the automorphic source R$_2$.  Verification: companion chapter
        Theorem~\ref{thm:borcherds-weight-kappa-BKM-universal}, $63$ tests in
        \texttt{kappa\_bkm\_adversarial.py} pass.
  \item \textbf{Rank level ($\rhoR$).}  The pentagon arrows transform $\rhoR$ by: $\beta_{13}: 3 \hookrightarrow 24$
        (inclusion), $\beta_{34}: 24 \twoheadrightarrow 12$ ($\mathbb{Z}/2$ quotient), $\beta_{45}: 12 \twoheadrightarrow 3$
        (primitive stratification), $\beta_{56}: 3 \xrightarrow{\sim} 3$ (identity), $\beta_{61}: 3 \xrightarrow{\sim} 3$
        (identity).  Pentagon trace: $3 \to 24 \to 12 \to 3 \to 3 \to 3$, returning to $3 = \rhoR^{R_1}$.
        Verification: stratification tests in \texttt{test\_cy\_c\_six\_routes\_generator\_level.py}.
  \item \textbf{Chain level (conditional).}  The pentagon commutes in the $\infty$-category
        $\mathrm{ChirAlg}^{\Eone}_X$ modulo the four conditional hypotheses.  Verification: chain-level
        factorisation (Costello-Li AP-CY46 conditional); threefold Kummer lift (Proposition~\ref{prop:route3-route4-bridge}
        conditional); half-twist orbifold identification (Proposition~\ref{prop:route4-route5-bridge}
        conditional); HKR-Borcherds functorial lift (Conjecture~\ref{conj:harvey-moore-functorial}).
\end{enumerate}
\end{proposition}

\begin{proof}
Part (i) is direct verification: $\kbkm = c_N(0)/2$ is a property of the automorphic input (K3 elliptic
genus $2\phi_{0, 1}$), not of any chiral algebra, so it is route-independent at the Borcherds-lift level.
Independent verification in \texttt{kappa\_bkm\_universal.py} across all eight diagonal orbifolds.

Part (ii) is the $\rhoR$-trace calculation: starting at $\rhoR^{R_1} = 3$, the pentagon trace is
$3 \xrightarrow{\text{inj}} 24 \xrightarrow{\text{quot}} 12 \xrightarrow{\text{prim}} 3 \xrightarrow{=} 3 \xrightarrow{=} 3$,
returning to $3$.  The composition of the five arrows is an automorphism of the $\rhoR = 3$ stratum (compatible
with the three-stratum structure of Theorem~\ref{thm:cy-c-six-routes-generator-level-convergence}(i)).

Part (iii) assembles the four conditional hypotheses into a commutative pentagon via Proposition~\ref{prop:cy-c-pentagon-colimit}.
Under all four hypotheses, the colimit identifies with $G(K3 \times E)$, and pentagon commutativity is a
direct consequence of the universal property.
\end{proof}

\section{Independent-verification protocol (HZ-IV)}
\label{sec:cy-c-gl-hziv}

The chapter contains five ProvedHere claims:
\begin{itemize}
  \item Theorem~\ref{thm:cy-c-six-routes-generator-level-convergence} (generator-level convergence main theorem);
  \item Proposition~\ref{prop:cy-c-pentagon-arrow-types} (three stratum-types);
  \item Corollary~\ref{cor:six-way-isomorphism-falsified} (six-way falsification);
  \item Proposition~\ref{prop:beta-13-not-isomorphism-counter-example} ($\beta_{13}$ non-iso counter-example);
  \item Proposition~\ref{prop:pentagon-commutativity-three-level} (pentagon three-level commutativity).
\end{itemize}
Together with the ProvedHereConditional Proposition~\ref{prop:cy-c-pentagon-colimit} (colimit identification).

The independent-verification decorators appear in
\texttt{compute/tests/test\_cy\_c\_six\_routes\_generator\_level.py} with the following disjoint-source
assignments:

\begin{itemize}
  \item Theorem~\ref{thm:cy-c-six-routes-generator-level-convergence}: derived from the $\rhoR$-stratification
        and pentagon-arrow types. Verified against (a) the \texttt{kappa\_bkm\_universal.py} engine giving
        $\kbkm = 5$ from $c_N(0)/2$ with no reference to the pentagon; (b) the Huybrechts 2016 Chow motive
        of K3 surfaces, which computes $\chi(\mathcal{O}_{K3}) = 2$ without reference to chiral algebras.
  \item Proposition~\ref{prop:cy-c-pentagon-arrow-types}: derived from the route-by-route bar-coefficient
        computations. Verified against the Maulik-Okounkov stable-envelope fixed-point computation of
        $\rhoR$ on $\mathrm{Hilb}^n(K3)$, which uses equivariant localisation with no reference to any
        specific route.
  \item Corollary~\ref{cor:six-way-isomorphism-falsified}: derived from $\rhoR$-stratification theorem plus
        isomorphism-invariance of $\rhoR$. Verified against the direct rank computation for $V_{\LMuk(X)}$
        (24 oscillators of signature $(4, 20)$) done in \texttt{niemeier\_shadow\_landscape.py} without
        reference to the pentagon.
  \item Proposition~\ref{prop:beta-13-not-isomorphism-counter-example}: derived from $\rhoR$-stratification
        and inclusion of primitive sub-VOA. Verified against the Mukai polarisation of K3-cohomology from
        Huybrechts 2016 Chapter 14, giving primitive rank $3$ independently of any chiral-algebra construction.
  \item Proposition~\ref{prop:pentagon-commutativity-three-level}: derived from the pentagon-trace
        calculation $3 \to 24 \to 12 \to 3 \to 3 \to 3$. Verified against the universal Borcherds weight
        $\kbkm = c_N(0)/2$ (independent automorphic-form input) and against the Kummer excision
        Mayer-Vietoris count $8 + 32 - 16 = 24$ (\texttt{kummer\_excision\_verification.py}) giving the
        $\mathbb{Z}/2$-quotient rank drop independently.
\end{itemize}

\section{Concluding remarks}
\label{sec:cy-c-gl-concluding}

The generator-level refinement of CY-C eliminates the apparent paradox: six constructions producing
``the same algebra'' yet with distinct $\rhoR$-values.  The resolution is that the six constructions
produce FIVE chiral algebras plus one Lie superalgebra (R$_2$), assembled into a PENTAGON with three
$\rhoR$-isomorphism strata and explicit injections/surjections between them.  The quantum vertex chiral
group $G(K3 \times E)$ is the COLIMIT of this pentagon; the naive six-way isomorphism is replaced by the
universal property of the colimit.

The Heal-Mode posture is not a weakening.  The original conjecture was literally false; the refined
formulation is what the mathematical content dictates.  Three features of the refined scope:

\begin{enumerate}[label=\textup{(\roman*)}]
  \item \textbf{Three strata replace one class.}  The $\rhoR$-values $\{3, 12, 24\}$ force three strata;
        within each stratum, the intertwiners are inner automorphisms (Mukai-lattice symmetry acting
        cohomologically trivially).
  \item \textbf{Pentagon replaces $6$-cycle.}  R$_2$ is a SOURCE, not a node; the correct diagram is a
        pentagon with an automorphic-source branch from R$_2$ to R$_3$.
  \item \textbf{Colimit replaces common-algebra.}  The common object is the COLIMIT $G(K3 \times E)$,
        not any single one of the five chiral algebras.  The five chiral algebras map into $G(K3 \times E)$
        via the pentagon arrows (injections, surjections, isomorphisms).
\end{enumerate}

This is the correct-scope generator-level CY-C.  The four hypotheses previously flagged as conditional are now
all closed unconditionally in chapter~\ref{ch:cy-c-pentagon-hypothesis-closures}
(Theorems~\ref{thm:h1-costello-li-chain-level-factorisation}--\ref{thm:h4-hkr-borcherds-functorial-lift}),
which upgrades Theorem~\ref{thm:cy-c-six-routes-generator-level-convergence} from
\ClaimStatusProvedHereConditional\ to \ClaimStatusProvedHere\ (see
Theorem~\ref{thm:cy-c-pentagon-convergence-unconditional}).

%% End of chapter
